\paragraph{\trTwistedComplex。}\label{sub:twisted-complexes-representatives}
设 $\scr{B}$ 是一个\trStrictlyUnital{}的 \trACategory，且其链层面上的 $\hom$
空间皆为有限维；并取 $\scr{A}=Tw\scr{B}$。在这一情形下，\trTwistFunctor{}具有
显式的代表元。具体地说，给定 $Y_0$、$Y_1\in\Ob\scr{A}$，取\trTensorProduct{}%
$\hom_{\scr{A}}(Y_0,Y_1)\otimes Y_0$，其中把第一个因子看作一个带%
\trBoundaryOperator{} $\partial=\mu^1_{\scr{A}}$ 的\trChainComplex，一如%
\cref{sub:basic-operations}中那样。按定义，
\begin{equation}\label{eq:tw-hom-dual-identification}
\hom_{\scr{A}}\bigl(\hom_{\scr{A}}(Y_0,Y_1)\otimes Y_0,\,Y_1\bigr)
=\hom_{\scr{A}}(Y_0,Y_1)^{\vee}\otimes\hom_{\scr{A}}(Y_0,Y_1).
\end{equation}
这个\trMorphism{}群上的 $\mu^1_{\scr{A}}$-算子是
$(c^{\vee}\otimes b)\mapsto(-1)^{|b|}\partial^{\vee}(c^{\vee})\otimes b
+c^{\vee}\otimes\partial(b)$，其中 $\partial^{\vee}$ 是通常意义下的对偶%
\trDifferential。特别地，若把\eqref{eq:tw-hom-dual-identification} 与
$\mathrm{End}_{\KK}(\hom_{\scr{A}}(Y_0,Y_1))$ 等同起来，使得 $c^{\vee}\otimes b$
变为\trEndomorphism{} $a\mapsto\langle c^{\vee},a\rangle b$，那么恒等%
\trEndomorphism{}就是一个\trCocycle。把\eqref{eq:tw-hom-dual-identification}
中相应的元素记作 $f$，
并按\cref{sub:mapping-cones}中的定义作出\trMappingCone{} $\mathrm{Cone}(f)$。
在\trYonedaEmbedding{}之下，$\hom_{\scr{A}}(Y_0,Y_1)\otimes Y_0$ 变成
$\hom_{\scr{A}}(Y_0,Y_1)\otimes\scr{Y}_0$；$\scr{J}^1(f)$ 就是
\eqref{eq:abstract-evaluation-homomorphism} 所定义的那个抽象%
\trEvaluationMap{} $\phi$；从而 $\scr{C}\mathit{one}(f)=\scr{T}_{Y_0}
(\scr{Y}_1)$。这意味着，在这一类\trACategory{} $\scr{A}$ 中，我们可以令
\begin{equation}\label{eq:twist-cone-representative}
T_{Y_0}Y_1=\mathrm{Cone}(f).
\end{equation}

\begin{remark}\label{thm:twist-explicit-representative}
把这件事再写得更显式一些或许不无教益，至少在 $Y_0$、$Y_1$ 是 $\scr{B}$ 中%
\trObject{}的这一简单情形下是如此。取 $\hom_{\scr{B}}(Y_0,Y_1)$ 的一组%
\trBasis{} $\{b^i\}$，以 $i\in I$ 标号，并记
$\mu^1_{\scr{B}}(b^i)=\sum_j\beta^{ji}b^j$。于是，$T_{Y_0}(Y_1)$ 就是如下的%
\trTwistedComplex：
\begin{equation}\label{eq:twist-explicit-object}
Y=\bigoplus_{i\in I}S^{-|b^i|+1}Y_1\oplus Y_0,
\end{equation}
其\trDifferential{} $\delta_Y=(\phi^{ji}\otimes a^{ji})$ 由诸 $\phi$ 组成，
它们在 $\KK$ 的各个\trSuspension{}副本之间充当恒等映射；而
\begin{equation}\label{eq:twist-explicit-differential}
a^{ji}=
\begin{cases}
(-1)^{\sigma^i+1}\beta^{ji}e_{Y_0} & \text{若 } i,j\in I,\\
(-1)^{\sigma^i}b^i & \text{若 } i\in I,\ j=\ast,\\
0, & \text{在其余情形。}
\end{cases}
\end{equation}
这里，$\ast$ 是\eqref{eq:twist-explicit-object} 中对应于 $Y_0$ %
\trDirectSummand{}的新增行与列的标号（另一种做法是，
若按\cref{thm:shifted-multiplicity-spaces}的方式更改
\eqref{eq:shifted-multiplicity-identification} 中的等同，则
\eqref{eq:twist-explicit-differential} 两行中的符号 $(-1)^{\sigma^i}$ 便会
消失，留下 $-\beta^{ji}e_{Y_0}$ 与 $b^i$）。
\end{remark}

\trCanonical{}代表元\eqref{eq:twist-cone-representative} 的存在性有许多
实际上的好处。
例如，\trAFunctor{} $T_{Y_0}$ 可以通过直接写下公式来定义，而不必像%
\cref{sub:twisting-yoneda-modules}中那样采用间接的构造。

\begin{remark}\label{thm:c-finite-reduction}
设 $\tilde{\scr{A}}$ 是一个 \trCfinite{}的\trTriangulated{}\trACategory。应用%
\trPerturbation{}引理，把它换成一个与之\trQuasiisomorphic{}的、满足
$\mu^1=0$ 的\trCategory——从而其链层面上的 $\hom$ 空间都是有限维的。接着，
再用\cref{thm:strictly-unital-modification}使该\trCategory{}成为%
\trStrictlyUnital{}的。把所得结果记作 $\scr{B}$，并令 $\scr{A}=Tw\scr{B}$。由
构造可知，$\scr{B}$ 与 $\tilde{\scr{A}}$ \trQuasiisomorphic，因而它也是一个%
\trTriangulated{}\trACategory{}，这又转而蕴含 $\scr{A}$ 与 $\tilde{\scr{A}}$
\trQuasiequivalence。利用这一点，人们常常可以把手头这个相当一般的%
\trCfinite{}\trTriangulated{}\trACategory{}情形，化归到我们刚才讨论的那个
技术上更简单的情形。
\end{remark}
